\clearpage
\begin{explanationbox}
\textbf{先用一条真实的弦认清所有字母。}图中 $y^2=4x$，故 $p=2$；直线以 $60^\circ$ 穿过焦点。
\begin{center}\begin{tikzpicture}[font=\small,>=stealth]
\draw[rounded corners=3pt,gray!45] (0,0) rectangle (14.2,12);
\node at (7.1,11.45) {先认图：焦点在弦内，两段焦半径相加才是整条弦};
\begin{scope}[shift={(3.8,5.6)},scale=1.1]
\draw[->,gray] (-1.5,0)--(4.5,0) node[above] {$x$};
\draw[teal!75!black,thick,samples=120,domain=-2.3:4.15,variable=\t] plot ({\t*\t/4},{\t});

\draw[violet,densely dashed] (-1,-2.4)--(-1,4.2);
\node[rotate=90,violet] at (-1.4,1.6) {准线 $x=-1$};
\draw[blue!65!black,dashed] (-1,3.464101615)--(3,3.464101615);
\draw[orange!85!black,dashed] (-1,-1.154700538)--(.333333333,-1.154700538);
\fill (-1,3.464101615) circle(1.5pt) node[left=4pt] {$M$};
\fill (-1,-1.154700538) circle(1.5pt) node[left=4pt] {$N$};
\draw[gray,dotted] (1,.5)--(1,4.15);
\draw[<->,violet] (-1,4.1)--(1,4.1) node[midway,above] {$p=2$};
\draw[blue!75!black,very thick] (1,0)--(3.000000000,3.464101615);
\draw[orange!85!black,very thick] (1,0)--(0.333333333,-1.154700538);
\fill (3.000000000,3.464101615) circle(2pt) node[above=8pt] {$A$};
\fill (0.333333333,-1.154700538) circle(2pt) node[below left=6pt] {$B$};
\fill (1,0) circle(2pt) node[below right=5pt] {$F$};

\node[below left=4pt] at (0,0) {$O$};
\draw[violet,->] (1.65,0) arc[start angle=0,end angle=60,radius=.65];
\node[violet] at (1.94,.42) {$\theta$};
\end{scope}
\node[anchor=west,align=left] at (9.3,7.2)
 {$y^2=4x$，$p=2$\\[9pt]$F=(1,0)$\\[12pt]
 $\textcolor{blue!75!black}{r=FA=AM=4}$\\[9pt]
 $\textcolor{orange!85!black}{s=FB=BN=4/3}$\\[12pt]
 $L=r+s=16/3$};
\node at (7.1,2.1) {定义：到焦点的距离 $=$ 到准线的垂直距离。};
\node at (7.1,1.3) {$A=(3,2\sqrt3),\quad B=(1/3,-2\sqrt3/3),\quad\theta=60^\circ$。};
\node at (7.1,.5) {$p$ 是焦点到准线的距离；方程系数与通径长度都是 $2p=4$。};
\end{tikzpicture}
\end{center}
蓝段是 $FA=r$，橙段是 $FB=s$。两端到准线的水平距离分别等于这两段长度，故
\[
 L=r+s=AM+BN=x_A+x_B+p.
\]
图中两段为 $4$ 与 $4/3$，所以整弦是 $16/3$，不能只算一段或任意倍增。
\end{explanationbox}
\clearpage
\begin{explanationbox}
\textbf{只改变方向，弦长怎样变化？}保持抛物线、焦点和绘图尺度相同，看两端的变化。
\begin{center}\begin{tikzpicture}[font=\small,>=stealth]
\draw[rounded corners=3pt,gray!45] (0,0) rectangle (14.2,11);
\node at (7.1,10.45) {只转动过焦点的直线：越接近竖直越短，互补角等长};
\node at (2.6,9.5) {$\theta=60^\circ$};
\begin{scope}[shift={(2.0,5.2)},scale=.68]
\draw[->,gray] (-1.5,0)--(4.4,0) node[above] {$x$};
\draw[teal!75!black,thick,samples=120,domain=-4:4,variable=\t] plot ({\t*\t/4},{\t});
\draw[blue!75!black,very thick] (1,0)--(3.000000000,3.464101615);
\draw[orange!85!black,very thick] (1,0)--(0.333333333,-1.154700538);
\fill (3.000000000,3.464101615) circle(2pt) node[above=8pt] {$A$};
\fill (0.333333333,-1.154700538) circle(2pt) node[below left=6pt] {$B$};
\fill (1,0) circle(2pt) node[below right=5pt] {$F$};
\end{scope}
\node at (2.6,1.65) {$L=16/3$};
\node at (7.0,9.5) {$\theta=90^\circ$};
\begin{scope}[shift={(6.4,5.2)},scale=.68]
\draw[->,gray] (-1.5,0)--(4.4,0) node[above] {$x$};
\draw[teal!75!black,thick,samples=120,domain=-4:4,variable=\t] plot ({\t*\t/4},{\t});
\draw[blue!75!black,very thick] (1,0)--(1.000000000,2.000000000);
\draw[orange!85!black,very thick] (1,0)--(1.000000000,-2.000000000);
\fill (1.000000000,2.000000000) circle(2pt) node[above=8pt] {$A$};
\fill (1.000000000,-2.000000000) circle(2pt) node[below left=6pt] {$B$};
\fill (1,0) circle(2pt) node[below right=5pt] {$F$};
\end{scope}
\node at (7.0,1.65) {$L=4$};
\node at (11.4,9.5) {$\theta=120^\circ$};
\begin{scope}[shift={(10.8,5.2)},scale=.68]
\draw[->,gray] (-1.5,0)--(4.4,0) node[above] {$x$};
\draw[teal!75!black,thick,samples=120,domain=-4:4,variable=\t] plot ({\t*\t/4},{\t});
\draw[blue!75!black,very thick] (1,0)--(0.333333333,1.154700538);
\draw[orange!85!black,very thick] (1,0)--(3.000000000,-3.464101615);
\fill (0.333333333,1.154700538) circle(2pt) node[above=8pt] {$A$};
\fill (3.000000000,-3.464101615) circle(2pt) node[below left=6pt] {$B$};
\fill (1,0) circle(2pt) node[above right=5pt] {$F$};
\end{scope}
\node at (11.4,1.65) {$L=16/3$};

\node at (7.1,.65) {三图均为 $y^2=4x$、$F=(1,0)$，同尺度；通径 $L_{\min}=2p=4$。};
\end{tikzpicture}
\end{center}
中间是竖直通径，两侧的焦半径相等，弦长最短；两边的弦关于 $x$ 轴对称，长短两段互换，总长相同。
\[
 L=\frac{2p}{\sin^2\theta}\ge2p,\qquad
 \sin^2(\pi-\theta)=\sin^2\theta.
\]
\textbf{读图要点：}越接近 $90^\circ$ 越短；不是角度越大就越短。
下面分别画出正向、反向的投影，解释公式中的两个分母从哪里来。
\end{explanationbox}
